\section{Derivation of Recursive Learning Models}\label{sec:appendix_learning}
\subsection{Recursive Least Squares}\label{app:RLS}

The first updating mechanism we consider is Recursive Least Squares (RLS), which is a recursive implementation of the standard Ordinary Least Squares (OLS).
In OLS, the agent estimates the coefficients of a linear regression model by minimizing the sum of squared residuals.
Assuming agents use OLS to estimate the parameters in the PLM (\ref{eq:plm_univariate}), the OLS estimator $\beta_t$ can be derived from the minimization\footnote{
We use the following timing:
at time $t$ the agent forecasts $r^e_{t+1} = x_t \beta_{t}$ which then implies a realization for $r_t$. 
Therefore, to avoid simultaneity issues, we assume that when computing the current estimate, the last available log-return observation is $r_{t-1}$.
} problem as
\begin{equation}
     \beta_t = \left(\sum^{t-1}_{i=1}  x_{i-1}^2 \right)^{-1} \sum^{t-1}_{i=1}  x_{i-1} r_i.
\end{equation}

Estimating OLS on the whole available data set is sometimes referred to as off-line estimation. 
Another approach is to estimate the model on-line, that is, as soon as a new data point arrives. 
It turns out that the estimator formula allows for a recursive formulation.
Define 
\begin{equation}
M_{t-1} = \frac{1}{t-1}  \left(\sum^{t-1}_{i=1}  x_{i-1}^2 \right), \quad \beta_t = \frac{1}{t-1} M^{-1}_{t-1} \sum^{t-1}_{i=1}  x_{i-1} r_i.
\end{equation}

First we establish a recursion for $M_t$. Start from 
\begin{equation}
\frac{t-1}{t} M_{t-1} = \frac{1}{t}  \left(\sum^{t-1}_{i=1}  x_{i-1}^2 \right) \implies \frac{t-1}{t} M_{t-1} + \frac{1}{t} x_{t-1}^2 = M_t,
\end{equation}
that is
\begin{equation}
    M_t = M_{t-1} + \frac{1}{t} \left(x_{t-1}^2 - M_{t-1} \right).
\end{equation}

For $\beta_t$ we have
\begin{equation}\label{eq:b_t_evolution}
    \begin{aligned}
        \beta_t & = \beta_{t-1} + \left[(t-1) M_{t-1} \right]^{-1} \left\{ x_{t-2} r_{t-1} + \left[(t-2) M_{t-2} \right] \beta_{t-1} \right\} - \beta_{t-1} \\
        & = \beta_{t-1} + \left[(t-1) M_{t-1} \right]^{-1} \left\{ x_{t-2} r_{t-1} - x_{t-2}^2 \beta_{t-1} \right\}.
    \end{aligned}
\end{equation}

\subsection{Constant Gain Learning}\label{app:WLS}
In Weighted Least Squares (WLS), the agent estimates the coefficients of a linear regression model by minimizing a weighted sum of squared residuals.
The weights are typically chosen to reflect the relative importance or reliability of different observations. Assuming the linear model (\ref{eq:plm_univariate}),
the WLS estimator can be derived from the minimization problem
\begin{equation}
    \min_{\beta} \sum^{t-1}_{i=1} {w_{t-1,i}} (r_i - \beta x_{i-1}) ^2,
\end{equation}
with solution
\begin{equation}
    {\beta} = \left(\sum^{t-1}_{i=1} {w_{t-1,i}} x_{i-1}^2 \right)^{-1} \sum^{t-1}_{i=1} {w_{t-1,i}} x_{i-1} r_i.
\end{equation}

In particular, we are interested in the case where the weights decline geometrically over time\footnote{
    In empirical applications, instead of using geometrically declining weights, it is common to use a rolling window approach, where only the most recent observations within a fixed-size window are used for estimation. 
    However, the latter is analytically intractable.
    It is possible to determine a relationship between the window size and the geometric decay parameter $\gamma$ such that the two methods give similar weights \emph{overall} on the $\ell$ most recent observations:
    $\ell \approx \log(1 - p) / \log(1 - \gamma)$, where $p$ is the proportion of total weight assigned to the $\ell$ most recent observations.
}
\begin{equation}
    w_{t,i} = \gamma (1 - \gamma)^{t-i}, \quad \gamma \in (0,1).
\end{equation}
Similar to the OLS case, a recursive relationship can be derived.

Define 
\begin{equation}
    M_{t-1} = \gamma \sum^{t-1}_{i=1} {(1 - \gamma)^{t-1-i}} x_{i-1}^2, \quad \beta_{t} = \gamma M^{-1}_{t-1} \sum^{t-1}_{i=1} {(1 - \gamma)^{t-1-i}} x_{i-1} r_i.
\end{equation}
First we establish a recursion for $M_t$. Start from
\begin{equation}
M_{t-1} = \gamma \sum^{t-1}_{i=1} {(1 - \gamma)^{t-1-i}} x_{i-1}^2,
\end{equation}
then
\begin{equation}
(1 - \gamma) M_{t-1} = \gamma \sum^{t-1}_{i=1} {(1 - \gamma)^{t-i}} x_{i-1}^2 = M_t - \gamma x_{t-1}^2.
\end{equation}
That is
\begin{equation}
    M_t = M_{t-1} + \gamma \left(x_{t-1}^2 - M_{t-1} \right).
\end{equation}

Similarly for $\beta_t$ we start from the lagged expression
\begin{equation}
    \beta_{t-1} = \gamma M^{-1}_{t-2} \sum^{t-2}_{i=1} {(1 - \gamma)^{t-2-i}} x_{i-1} r_i,
\end{equation}
then
\begin{equation}
    (1- \gamma) \beta_{t-1} = \gamma M^{-1}_{t-2} \sum^{t-1}_{i=1} {(1 - \gamma)^{t-1-i}} x_{i-1} r_i ,
\end{equation}
and
\begin{equation}
 \beta_{t-1} ( M_{t-1} - \gamma x_{t-2}^2) =  \gamma \sum^{t-1}_{i=1} {(1 - \gamma)^{t-1-i}} x_{i-1} r_i \rightarrow \beta_{t-1} (1 - \gamma M^{-1}_{t-1} x_{t-2}^2) = \gamma M^{-1}_{t-1} \sum^{t-1}_{i=1} {(1 - \gamma)^{t-1-i}} x_{i-1} r_i,
\end{equation}
since \[M_{t-2} = \frac{1}{1 - \gamma} (M_{t-1} - \gamma x_{t-2}^2).\]
Adding $\gamma M^{-1}_{t-1} x_{t-2} r_{t-1}$ to both sides we obtain
\begin{equation}
    \beta_{t-1} (1 - \gamma M^{-1}_{t-1} x_{t-2}^2) + \gamma M^{-1}_{t-1} x_{t-2} r_{t-1} = \beta_t,
\end{equation}
which can be rearranged as
\begin{equation}
    \beta_t = \beta_{t-1} + \gamma M^{-1}_{t-1} x_{t-2} (r_{t-1} - x_{t-2} \beta_{t-1}).
\end{equation}

\subsection{Lasso Learning}\label{app:lasso}
The LASSO (Least Absolute Shrinkage and Selection Operator) estimator is defined as the solution to the optimization problem
\[
\hat{\beta}_{\text{LASSO}} = \arg\min_{\beta} \left\{\frac{1}{2}\left[(Y - X\beta)'(Y - X\beta)\right] + \lambda \sum_{j=1}^{k} |\beta_j|\right\},
\]
where $\lambda \geq 0$ is a tuning parameter that controls the amount of regularization. The LASSO estimator can be interpreted as a trade-off between fitting the data well (minimizing the sum of squared residuals) and keeping the model simple (minimizing the absolute values of the coefficients).

In general there is no closed form solution for the LASSO estimator, since the absolute value function is not differentiable at zero. However, in the special case where $X$ is an orthonormal matrix, the LASSO estimator can be expressed in closed form.
To do so we can rewrite the optimization problem as
\[
\hat{\beta}_{\text{LASSO}} 
= \arg\min_{\beta} 
\Big\{ \frac{1}{2}\left[Y'Y - 2\beta'\underbrace{X'Y}_{\beta_{\text{OLS}}} 
+ \beta'\underbrace{X'X}_{I}\beta\right]
+ \lambda \sum_{j=1}^{k} |\beta_j| \Big\},
\]
where the underbrackets follow from orthonormality and notice that the term $Y'Y$ does not depend on $\beta$, so it does not change the minimizer.
Using $A'A = \sum_{j=1}^{k} a_j^2$, we can rewrite the problem as
\[
\hat{\beta}_{\text{LASSO}} = \arg\min_{\beta} \Big\{ \sum_{j=1}^{k} \left( -\beta_j \beta_{\text{OLS},j} + \frac{1}{2}\beta_j^2 + \lambda |\beta_j| \right) \Big\}.
\]

This problem is separable in the coefficients $\beta_j$, so we can solve for each coefficient independently. The optimization problem for each $j$ coefficient is
\[\min_{\beta_j} \left\{\frac{1}{2} \beta_j^2 - \beta_j \beta_{\text{OLS},j} + \lambda |\beta_j| \right\}.\]
To solve this problem, we consider two cases:
\begin{enumerate}
    \item $\beta_{\text{OLS},j} > 0$, then $\beta_j$ has to be positive, otherwise we could decrease the objective by setting $\beta_j = 0$. 
    In this case, the problem becomes 
    \[\min_{\beta_j \geq 0} \left\{ \frac{1}{2} \beta_j^2 - \beta_j \beta_{\text{OLS},j} + \lambda \beta_j \right\},\]
    and differentiating and setting to zero gives
    \[\beta_j - \beta_{\text{OLS},j} + \lambda = 0 \Rightarrow \beta_j = \beta_{\text{OLS},j} - \lambda.\]
    Imposing non-negativity yields
    \[\hat{\beta}_{\text{LASSO},j} = \max\left(0, \beta_{\text{OLS},j} - \lambda\right).\]
    \item $\beta_{\text{OLS},j} < 0$, then $\beta_j$ has to be negative, otherwise we could decrease the objective by setting $\beta_j = 0$. In this case, the problem becomes
    \[\min_{\beta_j < 0} \left\{ \frac{1}{2} \beta_j^2 - \beta_j \beta_{\text{OLS},j} + \lambda (-\beta_j) \right\},\]    
    and
    \[\beta_j - \beta_{\text{OLS},j} - \lambda = 0 \Rightarrow \beta_j = \beta_{\text{OLS},j} + \lambda.\]
    Imposing negativity yields
    \[\hat{\beta}_{\text{LASSO},j} = \min\left(0, \beta_{\text{OLS},j} + \lambda\right).\]
\end{enumerate}

Combing the two cases, we obtain 
\begin{equation}
\hat{\beta}_{\text{LASSO}} = S(\beta_{\text{OLS}}, \lambda) = \textrm{sign}(\beta_{\text{OLS}}) \cdot \max\left(0, |\beta_{\text{OLS}}| - \lambda\right),
\end{equation}
where $S(\cdot)$ is the soft-thresholding operator.


